\documentclass[11pt,a4paper]{article}
\usepackage[T1]{fontenc}
\usepackage[utf8]{inputenc}
\usepackage{lmodern}
\usepackage[english]{babel}
\usepackage[margin=2.6cm]{geometry}
\usepackage{microtype}
\usepackage{booktabs,longtable,array}
\usepackage{enumitem}
\usepackage{fancyhdr}
\usepackage{titlesec}
\usepackage[font=small,labelfont=bf,labelsep=period,skip=6pt]{caption}
\usepackage{xcolor}
\usepackage{needspace}
\usepackage[hidelinks]{hyperref}

\definecolor{accent}{HTML}{1F3A5F}
\definecolor{rowgray}{HTML}{F2F4F7}

\newcommand{\RP}{RacePulse}
\newcolumntype{L}[1]{>{\raggedright\arraybackslash}p{#1}}

\titleformat{\section}{\Large\bfseries\color{accent}}{\thesection}{0.8em}{}
\titleformat{\subsection}{\large\bfseries\color{accent}}{\thesubsection}{0.8em}{}
\titleformat{\subsubsection}{\normalsize\bfseries\color{accent}}{\thesubsubsection}{0.8em}{}
\titlespacing*{\section}{0pt}{0pt}{12pt}
\titlespacing*{\subsection}{0pt}{18pt}{8pt}
\titlespacing*{\subsubsection}{0pt}{14pt}{6pt}

\setlength{\parindent}{0pt}
\setlength{\parskip}{0.7em}
\renewcommand{\arraystretch}{1.25}
\setlength{\LTpre}{6pt}
\setlength{\LTpost}{6pt}

\pagestyle{fancy}
\fancyhf{}
\fancyhead[L]{\small\textcolor{accent}{\RP{} --- RASD}}
\fancyfoot[C]{\small\thepage}
\renewcommand{\headrulewidth}{0.4pt}
\renewcommand{\headrule}{\hbox to\headwidth{\color{accent}\leaders\hrule height \headrulewidth\hfill}}

\begin{document}

% ---------------------------------------------------------------- Front page
\begin{titlepage}
\thispagestyle{empty}
\centering
\vspace*{2.2cm}
{\large\color{accent}\scshape Politecnico di Milano\par}
{\normalsize Software Engineering 2 \quad|\quad A.Y.\ 2026--2027\par}
\vspace{2.6cm}
{\fontsize{34}{40}\selectfont\bfseries\color{accent}\RP\par}
\vspace{0.6cm}
{\Large Requirements Analysis and\\[2pt] Specification Document\par}
\vspace{0.4cm}
{\large (RASD)\par}
\vspace{1.4cm}
\textcolor{accent}{\rule{6cm}{0.8pt}}\par
\vspace{1.4cm}
{\large\bfseries Authors\par}
\vspace{0.4cm}
{\large Giulia Blanchetti \quad 10913272\par}
{\large Nicola Cornali \quad 10919230\par}
\vfill
\begin{tabular}{rl}
\end{tabular}
\vspace{1.5cm}
\end{titlepage}

% ---------------------------------------------------------------- Contents
\pagenumbering{roman}
\setcounter{page}{2}
\renewcommand{\contentsname}{Table of Contents}
\tableofcontents
\clearpage

\pagenumbering{arabic}

% ---------------------------------------------------------------- Chapter 1
\section{Introduction}

\subsection{Purpose}
Formula~1 racing sessions generate a large, continuous stream of heterogeneous data, including car telemetry,
real-time positions, lap and sector times, driver intervals, pit-stop details, weather conditions, and
race-control events. Different categories of consumers require access to distinct subsets of this information:
racing enthusiasts seek to follow or replay races through visual interfaces, while developers of third-party
software systems require programmatic access to analyze telemetry, compare performances across sessions, or
track driver progression.

\RP{} is a distributed software platform designed to ingest this data from an external provider (OpenF1 in our
case), store it persistently, and distribute it---either in near real time or on demand---to a large and
heterogeneous population of consumers via dedicated graphical and application programming interfaces (GUI and
API). The platform ensures customized delivery, high availability, and timeliness without suffering uncontrolled
data loss during workload peaks or propagating consumer-side failures.

The goals of \RP{} are listed in Table~\ref{tab:goals}.

\begin{longtable}{@{}L{1.4cm} L{13.3cm}@{}}
\caption{Goals of \RP.}\label{tab:goals}\\
\toprule \textbf{ID} & \textbf{Goal}\\ \midrule \endfirsthead
\toprule \textbf{ID} & \textbf{Goal}\\ \midrule \endhead
\bottomrule \endfoot
G1 & The system continuously acquires Formula~1 racing session information and telemetry from the external data provider.\\
G2 & A consumer can follow an ongoing racing session in near real time.\\
G3 & A consumer can customize and filter the specific categories of racing data and events they want to receive.\\
G4 & A consumer can query, analyze, and replay historical data from past racing sessions.\\
G5 & Third-party software systems can consume live and historical racing data programmatically.\\
G6 & The system persistently stores all acquired racing session data to guarantee long-term availability.\\
G7 & A consumer can access the platform without their service availability or performance being degraded by the failures or abnormal behavior of other consumers.\\
\end{longtable}

\subsection{Scope}
The scope of \RP{} is defined by the boundary between the machine (the software system to be built) and the
world (the operational environment with which it interacts). This section delineates the system perimeter,
formalizes the atomic phenomena crossing the boundary, and establishes clear interpretations for the ambiguous
aspects of the problem statement.

\subsubsection{World phenomena}
Table~\ref{tab:wp} lists the phenomena that take place in the world and are not observed or controlled by \RP{}.
 
\begin{longtable}{@{}L{1.4cm} L{13.3cm}@{}}
\caption{World phenomena.}\label{tab:wp}\\
\toprule \textbf{ID} & \textbf{Description}\\ \midrule \endfirsthead
\toprule \textbf{ID} & \textbf{Description}\\ \midrule \endhead
\bottomrule \endfoot
WP1 & A Formula~1 racing session starts, evolves, and ends.\\
WP2 & Weather conditions change at the circuit.\\
WP3 & Race-control events (e.g., red flags, safety car interventions) occur during the race.\\
WP4 & OpenF1 collects car telemetry, including speed, throttle, braking, gear, and RPM.\\
WP5 & OpenF1 collects approximate car positions on the circuit.\\
WP6 & OpenF1 collects lap and sector times, driver positions, and intervals (gaps).\\
WP7 & OpenF1 collects pit-stop information.\\
WP8 & OpenF1 collects weather condition updates from the circuit.\\
WP9 & OpenF1 collects race-control events generated by the race direction.\\
\end{longtable}
 
\subsubsection{Shared phenomena}
Table~\ref{tab:sp} lists the phenomena shared between the world and the machine, with the entity that controls
each of them and the entity that observes it.
 
\begin{longtable}{@{}L{1.2cm} L{8.9cm} L{2.1cm} L{2.1cm}@{}}
\caption{Shared phenomena.}\label{tab:sp}\\
\toprule \textbf{ID} & \textbf{Description} & \textbf{Controller} & \textbf{Observer}\\ \midrule \endfirsthead
\toprule \textbf{ID} & \textbf{Description} & \textbf{Controller} & \textbf{Observer}\\ \midrule \endhead
\bottomrule \endfoot
SP1 & \RP{} requests the live data stream for an ongoing session from OpenF1. & \RP & OpenF1\\
SP2 & \RP{} requests historical race data from OpenF1 to populate or sync its persistent storage. & \RP & OpenF1\\
\midrule
SP3 & OpenF1 sends the requested car telemetry (speed, throttle, braking, gear, RPM) to \RP. & OpenF1 & \RP\\
SP4 & OpenF1 sends the requested approximate car positions to \RP. & OpenF1 & \RP\\
SP5 & OpenF1 sends lap times, sector times, and driver intervals to \RP. & OpenF1 & \RP\\
SP6 & OpenF1 sends pit-stop information to \RP. & OpenF1 & \RP\\
SP7 & OpenF1 sends weather condition updates to \RP. & OpenF1 & \RP\\
SP8 & OpenF1 sends race-control events (flags, safety car announcements) to \RP. & OpenF1 & \RP\\
\midrule
SP9 & The Consumer connects and authenticates to the \RP{} platform via GUI or API. & Consumer & \RP\\
SP10 & The Consumer selects an ongoing racing session to monitor. & Consumer & \RP\\
SP11 & The Consumer specifies the data categories they want to receive (e.g., telemetry, weather, session info). & Consumer & \RP\\
SP12 & The Consumer specifies the drivers they want to monitor. & Consumer & \RP\\
SP13 & The Consumer modifies their active data or driver subscription choices during the session. & Consumer & \RP\\
SP14 & The Consumer submits a query for historical data, specifying sessions, drivers, laps, weather, or telemetry. & Consumer & \RP\\
SP15 & The Consumer requests the replay of a specific past session. & Consumer & \RP\\
\midrule
SP16 & \RP{} delivers the tailored live data stream to the Consumer based on their active subscriptions. & \RP & Consumer\\
SP17 & \RP{} delivers time-critical race-control alerts (e.g., red flags) to the subscribed Consumer. & \RP & Consumer\\
SP18 & \RP{} delivers the requested historical data logs to the Consumer. & \RP & Consumer\\
SP19 & \RP{} delivers the replay data stream, reproducing the race evolution exactly as it happened. & \RP & Consumer\\
\end{longtable}
 
\subsection{Definitions, Acronyms, Abbreviations}
 
\subsubsection{Definitions}
\begin{longtable}{@{}L{3.6cm} L{11.1cm}@{}}
\caption{Definitions.}\label{tab:defs}\\
\toprule \textbf{Term} & \textbf{Definition}\\ \midrule \endfirsthead
\toprule \textbf{Term} & \textbf{Definition}\\ \midrule \endhead
\bottomrule \endfoot
Consumer & A human user, typically a racing enthusiast, who obtains information from \RP{} through the GUI.\\
Third-party system & An external software system, developed by third-party developers, that obtains information from \RP{} programmatically through the API. It is a different entity from the Consumer.\\
Racing session & A single on-track event of a Formula~1 race weekend (e.g., practice, qualifying, sprint, race), identified by the external provider.\\
Live data & Information about a racing session that is still in progress, as received by \RP{} from the external provider.\\
Historical data & Information about racing sessions that have ended and that is stored persistently by \RP.\\
Data category & A kind of racing information handled by \RP: car telemetry, car positions, lap and sector times, driver intervals, pit-stop information, weather conditions, race-control events.\\
Race-control event & An event generated by the race direction during a session, such as a flag or a safety-car intervention.\\
Subscription & The choice made by a Consumer or a Third-party system of the session, data categories and drivers it wants to receive; it can be modified at any time.\\
Replay & The delivery to a Consumer or a Third-party system of the data of a past racing session, reproducing the evolution of the race as it happened.\\
Near real time & Under normal operating conditions, at least 95\% of the information received by \RP{} is made available to interested Consumers and Third-party systems within 2 seconds of its reception, and particularly important race-control events (e.g., red flags, safety-car announcements) within 1 second (assignment, Section~6.3).\\
Uncontrolled data loss & The silent loss of information that Consumers and Third-party systems reasonably expect to receive. If \RP{} cannot process incoming information immediately, a short delay is acceptable, whereas silent loss is not.\\
External data provider & OpenF1, the source of all racing information used by \RP.\\
\end{longtable}
 
\subsubsection{Acronyms and abbreviations}
\begin{longtable}{@{}L{3.6cm} L{11.1cm}@{}}
\caption{Acronyms and abbreviations.}\label{tab:acr}\\
\toprule \textbf{Acronym} & \textbf{Meaning}\\ \midrule \endfirsthead
\toprule \textbf{Acronym} & \textbf{Meaning}\\ \midrule \endhead
\bottomrule \endfoot
API & Application Programming Interface\\
DD & Design Document\\
F1 & Formula~1\\
G & Goal\\
GUI & Graphical User Interface\\
RASD & Requirements Analysis and Specification Document\\
RPM & Revolutions Per Minute\\
SP & Shared Phenomenon\\
WP & World Phenomenon\\
\end{longtable}
 
\Needspace{8\baselineskip}
\subsection{Revision history}
\begin{longtable}{@{}L{2cm} L{3.2cm} L{9.5cm}@{}}
\caption{Revision history.}\label{tab:rev}\\
\toprule \textbf{Version} & \textbf{Date} & \textbf{Description}\\ \midrule \endfirsthead
\toprule \textbf{Version} & \textbf{Date} & \textbf{Description}\\ \midrule \endhead
\bottomrule \endfoot
&&\\
\end{longtable}
 
\subsection{Reference documents}
\begin{itemize}[leftmargin=*,itemsep=3pt]
\item Politecnico di Milano, Software Engineering 2, A.Y.\ 2026--2027. \emph{Requirement Engineering and Design Project: goal, schedule, and rules} (assignment document).
\item M.\ Camilli, E.\ Di Nitto, M.\ G.\ Rossi. \emph{Software Engineering 2: lecture slides} (Requirements Engineering and Structure of a RASD document).
\item M.\ Camilli, E.\ Di Nitto, A.\ Fuggetta, A.\ Margara, M.\ G.\ Rossi, D.\ A.\ Tamburri. \emph{Software Engineering: a Structured Journey from Inception to Delivery}. BUP, 2025.
\item ISO/IEC/IEEE 29148:2018. \emph{Systems and software engineering --- Life cycle processes --- Requirements engineering}. \url{https://doi.org/10.1109/IEEESTD.2018.8559686}
\item M.\ Jackson, P.\ Zave. \emph{Deriving Specifications from Requirements: an Example}. 17th International Conference on Software Engineering, 1995.
\item OpenF1, open-source provider of historical and live Formula~1 data. \url{https://openf1.org/}
\end{itemize}
 
\subsection{Document structure}
This document follows the structure suggested by the course, derived from the IEEE standard, and is organized in six chapters.
\begin{description}[leftmargin=!,labelwidth=3.4cm,itemsep=4pt]
\item[Chapter 1: Introduction] presents the purpose of \RP{} and its goals, the scope of the analysis with the world and shared phenomena, and the definitions, revision history and reference documents.
\item[Chapter 2: Overall description] describes the product perspective (scenarios, shared phenomena and domain model), the main product functions, the characteristics of the users, and the domain assumptions, dependencies and constraints.
\item[Chapter 3: Specific requirements] details the external interface requirements, the functional requirements (use cases and their mapping on the requirements), the performance requirements, the design constraints and the software system attributes.
\item[Chapter 4: Formal analysis using Alloy] presents the objectives of the formal modeling activity, the Alloy model, and the results of its validation.
\item[Chapter 5: Effort spent] reports the hours worked by each group member on this document.
\item[Chapter 6: References] lists the sources used in this document.
\end{description}
\end{document}
